\documentclass[conference]{IEEEtran}
\IEEEoverridecommandlockouts

\usepackage{cite}
\usepackage{amsmath,amssymb,amsfonts}
\usepackage{algorithmic}
\usepackage{algorithm}
\usepackage{graphicx}
\usepackage{textcomp}
\usepackage{xcolor}
\usepackage{booktabs}
\usepackage{url}
\usepackage{hyperref}
\usepackage{multirow}
\usepackage{array}

\hypersetup{
    colorlinks=true,
    linkcolor=blue,
    citecolor=blue,
    urlcolor=blue
}

\def\BibTeX{{\rm B\kern-.05em{\sc i\kern-.025em b}\kern-.08em
    T\kern-.1667em\lower.7ex\hbox{E}\kern-.125emX}}

% ─── custom commands ──────────────────────────────────────────────────────────
\newcommand{\udop}{\textsc{udop}}
\newcommand{\docvqa}{DocVQA}
\newcommand{\anls}{ANLS}

\begin{document}

% ─── TITLE ───────────────────────────────────────────────────────────────────
\title{DocChat: Advanced Visual Document Intelligence and\\
Information Extraction using Multimodal Transformers}

% ─── AUTHORS ─────────────────────────────────────────────────────────────────
\author{
\IEEEauthorblockN{Rayan Ahmed}
\IEEEauthorblockA{
    \textit{Dept. of Artificial Intelligence} \\
    \textit{FAST School of Computing, NUCES} \\
    Islamabad, Pakistan \\
    i230018@isb.nu.edu.pk}
\and
\IEEEauthorblockN{Raza Ahmad}
\IEEEauthorblockA{
    \textit{Dept. of Artificial Intelligence} \\
    \textit{FAST School of Computing, NUCES} \\
    Islamabad, Pakistan \\
    i230060@isb.nu.edu.pk}
\and
\IEEEauthorblockN{Awwab Ahmad}
\IEEEauthorblockA{
    \textit{Dept. of Artificial Intelligence} \\
    \textit{FAST School of Computing, NUCES} \\
    Islamabad, Pakistan \\
    i230079@isb.nu.edu.pk}
}

\maketitle

% ─── ABSTRACT ────────────────────────────────────────────────────────────────
\begin{abstract}
Visual document understanding demands joint reasoning over textual content,
spatial layout, and visual rendering---modalities that standard language models
fundamentally cannot exploit after one-dimensional OCR linearization discards
all positional context. This paper presents \textbf{DocChat}, a fully
operational Visual Document Question Answering (DocVQA) system built around
Microsoft's Unified Document Processing (\udop) transformer
(\texttt{microsoft/udop-large-512-300k}), a multimodal encoder-decoder that
jointly encodes text tokens, normalized 2D bounding-box coordinates, and
patch-level visual embeddings within a single Vision-Text-Layout (VTL)
encoder. DocChat exposes a FastAPI backend executing PyTorch inference and a
Flutter Web dashboard supporting three analysis strategies: Single-Page, Smart
(TF-weighted OCR page ranking), and exhaustive All-Pages. To bridge the latency
gap between heavy academic transformers and real-time applications, we design a
two-level in-memory caching architecture. Level-1 caching provides instant
retrieval (0.0\,s) for repeated queries via MD5 document fingerprinting;
Level-2 caching reuses pre-rendered PIL page images and pre-computed Tesseract
OCR layouts, yielding a \textbf{5.6$\times$ speedup} (${\approx}$3.8\,s
vs.\ ${\approx}$21.4\,s cold-start) for subsequent distinct queries.
A Least Recently Used (LRU) eviction policy bounds memory to five concurrent
sessions. Evaluated on a custom three-document benchmark---invoice, receipt, and
student registration form---with standardized JSON annotations, DocChat achieves
\textbf{100\% extraction precision} on all 18 annotated query-answer pairs,
including structured financial entities such as invoice totals and alphanumeric
identifiers. The complete source code, evaluation dataset, and setup guides are
publicly open-sourced at \url{https://github.com/Rayan581/visual-doc-qa}.
\end{abstract}

\begin{IEEEkeywords}
visual document understanding, document VQA, multimodal transformers, UDOP,
layout-aware NLP, OCR, intelligent caching, FastAPI, Flutter
\end{IEEEkeywords}

% ═════════════════════════════════════════════════════════════════════════════
\section{Introduction}
\label{sec:intro}
% ═════════════════════════════════════════════════════════════════════════════

The digitization of administrative, legal, and financial workflows has produced
vast corpora of visually rich documents---invoices, receipts, regulatory forms,
contracts---whose semantic content is encoded not merely in character sequences
but in the interplay of spatial layout coordinates, typographic hierarchy, and
tabular grid structure. Extracting information from such documents is a critical
enterprise need, yet it remains unsolved by classical text-only pipelines.

Consider a standard financial invoice: the string \texttt{\$205.70} acquires
its meaning as \emph{Total Amount Due} not from token context alone, but from
its 2D alignment beneath a column header. A language model receiving a
linearized OCR transcript---stripped of all spatial metadata---faces an
ambiguous token stream in which six dollar-valued fields are
indistinguishable by position. This is not a capacity failure; it is a
representation failure.

Visual Document Question Answering (\docvqa) formalizes the task of answering
natural-language queries about document images by jointly reasoning over text,
layout, and visual appearance \cite{b_mathew2021}. Motivated by this challenge,
we present \textbf{DocChat}, which makes the following contributions:

\begin{enumerate}
    \item A production-grade \docvqa{} pipeline deploying \udop---the current
          state-of-the-art unified multimodal document model \cite{b_tang2023}---
          within a real-time FastAPI/Flutter system architecture.

    \item A \textbf{dual-level in-memory caching architecture} that reduces
          cold-start latency (${\approx}$21.4\,s) to ${\approx}$3.8\,s (Level-2
          hit) and 0.0\,s (Level-1 hit), enabling interactive document Q\&A.

    \item A \textbf{TF-weighted Smart Page Ranking} algorithm that achieves
          100\% extraction accuracy at roughly half the inference cost of
          exhaustive page traversal.

    \item A custom three-document evaluation benchmark with standardized
          JSON ground-truth annotations covering invoices, receipts, and
          administrative forms.
\end{enumerate}

To promote reproducibility and open-source scientific replication, all project source code, Flutter deployment modules, and the curated evaluation dataset are publicly hosted at \url{https://github.com/Rayan581/visual-doc-qa}.

The remainder of this paper is organized as follows.
Section~\ref{sec:related} reviews related work.
Section~\ref{sec:methodology} details the system architecture and algorithms.
Section~\ref{sec:dataset} describes the evaluation dataset.
Section~\ref{sec:experiments} reports experimental results.
Section~\ref{sec:discussion} discusses findings and limitations.
Section~\ref{sec:conclusion} concludes.

% ═════════════════════════════════════════════════════════════════════════════
\section{Related Work}
\label{sec:related}
% ═════════════════════════════════════════════════════════════════════════════

\subsection{Traditional OCR-Based Pipelines}

Early document information extraction systems decomposed the problem into two
independent stages: OCR followed by rule-based post-processing using regular
expressions or template matching \cite{b_smith2007}. While effective on fixed
templates, these systems exhibited brittleness to layout variation and could
not support free-form natural language queries. Hybrid Tesseract + BERT
baselines reported by the DocVQA challenge organizers achieved F1 scores of
56--67\%, compared to near-98\% human performance \cite{b_mathew2021}, precisely
because BERT receives the linearized OCR stream without spatial coordinates.

\subsection{The LayoutLM Family}

\textbf{LayoutLM} \cite{b_xu2020} was the first model to augment BERT token
embeddings with 2D position embeddings from OCR bounding boxes
$(x_0, y_0, x_1, y_1)$, pre-trained on 11M IIT-CDIP documents. This enabled
structural row/column relationship modeling regardless of token-sequence
distance. \textbf{LayoutLMv2} \cite{b_xu2021} incorporated a ResNeXt-101
visual backbone and a Text-Image Alignment pre-training objective, while
\textbf{LayoutLMv3} \cite{b_huang2022} introduced a symmetric text-image
encoder with Masked Language Modeling, Masked Image Modeling, and Word-Patch
Alignment objectives, eliminating the separate CNN encoder. However, all three
models retain an encoder-only BERT heritage, constraining output to extractive
spans rather than generative natural language answers, and all are limited to
512 input tokens.

\subsection{UDOP: Unified Document Processing}

\udop{} \cite{b_tang2023} resolves both limitations. Built upon the T5
encoder-decoder framework, it frames \emph{all} document intelligence tasks as
text-to-text generation, enabling generative answer production. Its
Vision-Text-Layout (VTL) unified encoder processes text tokens and image
patches within a single transformer stack from layer one, with shared relative
2D positional encodings enabling deep cross-modal attention. Pre-trained on 11M
documents from IIT-CDIP and CC-PDF for 300k steps, \udop-Large achieves 84.7\%
ANLS on the DocVQA leaderboard, outperforming LayoutLMv3 (83.4\%) and all
encoder-only baselines.

\subsection{OCR-Free Architectures}

Donut \cite{b_kim2022} and Nougat \cite{b_blecher2023} demonstrate that
end-to-end image-to-text transformers can achieve competitive document
understanding without any intermediate OCR stage, eliminating the
Tesseract dependency entirely. These represent the frontier for future
iteration of DocChat.

% ═════════════════════════════════════════════════════════════════════════════
\section{System Architecture and Methodology}
\label{sec:methodology}
% ═════════════════════════════════════════════════════════════════════════════

\subsection{End-to-End Pipeline Overview}
\label{subsec:pipeline}

Fig.~\ref{fig:pipeline} illustrates DocChat's eight-stage processing pipeline.
A document upload (PDF or image) and a natural-language query $q$ enter the
FastAPI backend. The system first computes an MD5 fingerprint of the raw file
bytes and consults the two-level cache before performing any expensive
computation. On a full cache miss, the pipeline executes document parsing,
Tesseract OCR, optional Smart page ranking, and finally UDOP inference,
caching intermediate products at each stage for future reuse.

\begin{figure}[t]
\centering
\fbox{%
\parbox{0.9\columnwidth}{%
\small
\textbf{Upload} (PDF/Image, query $q$)
$\;\rightarrow\;$ \textbf{MD5 Hash}
$\;\rightarrow\;$ \textbf{L1 Cache?} $\xrightarrow{\text{hit}}$ Return\\[2pt]
$\downarrow$ miss\\[2pt]
\textbf{L2 Cache?} $\xrightarrow{\text{hit}}$ Load PIL + OCR\\
$\downarrow$ miss\\[2pt]
\textbf{PDF/Image Parse} (pdf2image, 200\,DPI)\\
$\;\rightarrow\;$ \textbf{Tesseract OCR} (--psm\,11, normalize $[0,1000]$)\\
$\;\rightarrow\;$ \textbf{Page Ranking} (Smart) or All/Single\\
$\;\rightarrow\;$ \textbf{UDOP Inference} (CUDA/CPU)\\
$\;\rightarrow\;$ \textbf{Store L1 \& L2} $\;\rightarrow\;$ \textbf{Return Answer}
}}
\caption{DocChat end-to-end processing pipeline with dual-level cache routing.}
\label{fig:pipeline}
\end{figure}

\subsection{UDOP Inference Stage}

For each selected page, the \texttt{UdopProcessor} tokenizes the question
$q$ together with the OCR word sequence $\mathcal{W} = \{w_1, \ldots, w_N\}$
and their normalized bounding boxes
$\mathcal{B} = \{b_i = (x_0^i, y_0^i, x_1^i, y_1^i)\}_{i=1}^N$,
where all coordinates are integers in $[0, 1000]$ relative to page dimensions.

Each token $w_i$ receives a trimodal input embedding:
\begin{equation}
    \mathbf{e}_i = \mathbf{e}^{\text{tok}}_i
                 + \mathbf{e}^{\text{spatial}}_i
                 + \mathbf{e}^{\text{visual}}_i,
    \label{eq:embed}
\end{equation}
where the spatial embedding is:
\begin{align}
    \mathbf{e}^{\text{spatial}}_i &=
        E_x(x_0^i) + E_x(x_1^i)
      + E_y(y_0^i) + E_y(y_1^i) \notag \\
      &\quad + E_h(y_1^i - y_0^i)
              + E_w(x_1^i - x_0^i),
    \label{eq:spatial}
\end{align}
and $\mathbf{e}^{\text{visual}}_i$ is the patch embedding of the image region
corresponding to $b_i$. The encoder processes all modalities jointly through
every self-attention layer; the autoregressive decoder then generates the
answer string token-by-token.

\subsection{Smart Page Ranking: TF-Weighted Scoring}
\label{subsec:smart}

For multi-page documents, exhaustive per-page UDOP inference is
computationally prohibitive. The Smart strategy implements a lightweight
pre-filtering stage. Let $Q$ denote the set of normalized (lowercased,
stop-word-filtered) question tokens and $P_k$ the OCR token \emph{multiset}
for page $k$. Page relevance is scored as:
\begin{equation}
    \text{Score}(P_k,\, Q)
    = \sum_{q\,\in\, Q \cap P_k} \bigl(1 + \ln\, \text{TF}(q, P_k)\bigr),
    \label{eq:score}
\end{equation}
where $\text{TF}(q, P_k)$ is the raw term count of token $q$ in page $k$'s
OCR output. The logarithmic dampening prevents high-frequency tokens from
dominating: a token occurring once contributes weight 1.0; occurring
$e{\approx}2.718$ times contributes 2.0; occurring 100 times contributes
${\approx}5.6$. Pages are ranked in descending order by
\eqref{eq:score} and only the top-$N$ (default $N{=}2$) proceed to UDOP
inference.

\begin{algorithm}[t]
\caption{Smart Page Ranking}
\label{alg:smart}
\begin{algorithmic}[1]
\REQUIRE Question $q$, OCR multisets $\{P_k\}_{k=1}^K$, threshold $N$
\STATE $Q \leftarrow \text{tokenize\_normalize}(q)$
\FOR{$k = 1$ \TO $K$}
    \STATE $s_k \leftarrow 0$
    \FOR{each $w \in Q \cap P_k$}
        \STATE $s_k \mathrel{+}= 1 + \ln(\text{TF}(w, P_k))$
    \ENDFOR
\ENDFOR
\STATE $\mathcal{S} \leftarrow \text{argsort}(\{s_k\}, \text{descending})$
\RETURN $\mathcal{S}_{1:N}$
\end{algorithmic}
\end{algorithm}

\subsection{Dual-Level In-Memory Caching}
\label{subsec:cache}

\subsubsection{Level-1 Cache: Q\&A Fingerprinting}

The Level-1 cache maps composite document-question fingerprints to precomputed
answer strings. The key is:
\begin{equation}
    k_{\text{L1}} = \text{MD5}(B) \;\|\; \text{MD5}(\text{norm}(q)),
    \label{eq:l1key}
\end{equation}
where $B$ denotes the raw uploaded file bytes and $\|$ denotes string
concatenation. An MD5 key collision is cryptographically negligible for this
use case. On a Level-1 hit, the stored answer is returned in $O(1)$ with
zero computation; measured latency is 0.0\,s.

\subsubsection{Level-2 Cache: Document Pre-Processing Reuse}

The Level-2 cache stores the two most expensive intermediate products, keyed by
$\text{MD5}(B)$ alone:
\begin{itemize}
    \item \textbf{Rendered PIL images}: All pages at 200\,DPI (eliminating
          ${\approx}$4.2\,s of pdf2image/Poppler decoding per 3-page document).
    \item \textbf{Tesseract OCR layouts}: Per-page dictionaries of word strings
          and normalized bounding boxes (eliminating ${\approx}$5.1\,s/page
          of CPU-bound OCR computation).
\end{itemize}
On a Level-2 hit, only UDOP tokenization and inference remain
(${\approx}$3.8\,s), yielding the observed $5.6\times$ speedup.

\subsubsection{LRU Eviction}

Both caches are managed by an LRU policy
(\texttt{collections.OrderedDict}), retaining a maximum of $M{=}5$ active
sessions. The peak RAM footprint is bounded by:
\begin{equation}
    \text{RAM}_{\max} \approx M \cdot N_p \cdot W \cdot H \cdot 3,
    \label{eq:ram}
\end{equation}
where $N_p$ is the maximum page count and $W{\times}H$ is the rendered
page resolution in pixels. For $M{=}5$, $N_p{=}10$, and A4 at 200\,DPI
($1654{\times}2339$), this yields ${\approx}1.7$\,GB of PIL image
storage---acceptable for an 8\,GB server. OCR metadata contributes
negligibly ($<$5\,MB per session).

\subsection{Frontend: Flutter Web Dashboard}

The client is a Flutter Web application providing: a drag-and-drop file picker
with animated upload progress; a real-time chat stream rendering Q\&A pairs
conversationally; analysis strategy selector chips (\emph{Single Page} /
\emph{Smart} / \emph{All Pages}); metadata tag chips displaying cache status
(Level-1 Hit / Level-2 Hit / Cold Start) and latency in milliseconds. All
communication is via REST over CORS-enabled FastAPI endpoints, with documents
streamed as \texttt{multipart/form-data} and answers returned as JSON.

% ═════════════════════════════════════════════════════════════════════════════
\section{Dataset}
\label{sec:dataset}
% ═════════════════════════════════════════════════════════════════════════════

No publicly available benchmark precisely targets the combination of financial
invoices, retail receipts, and institutional registration forms most prevalent
in enterprise extraction workflows. We therefore constructed the
\textbf{DocChat Evaluation Benchmark}, comprising three annotated document
images and a standardized \texttt{annotations.json} ground-truth file.

\subsection{Document Corpus}

\textbf{sample\_invoice.png} is a financial Accounts Payable invoice
containing: vendor/purchaser header blocks, a ruled tabular line-items section
(description, quantity, unit price, extended amount), subtotal and GST rows, a
\emph{Total Amount Due} field (\$205.70), and a unique invoice identifier
(GT-98432). The tabular column-header-to-value association constitutes the
primary layout reasoning challenge.

\textbf{sample\_receipt.png} is a scanned point-of-sale receipt exhibiting
narrow-column thermal-printer layout: item names left-aligned, prices
right-aligned, quantity multiplier notation (\texttt{3x \$4.50}), explicit
GST rate, and a total payment row. The compressed vertical structure creates
OCR tokenization challenges where tokens from distinct semantic fields become
interleaved in linearized reading order.

\textbf{sample\_form.png} is a FAST-NUCES Student Registration Form with
labeled text fields (full name, student ID, program, semester, registration
date), section dividers, and checkboxes. This document tests field-label-to-value
association where labels are typographically bold and left-aligned.

\subsection{Annotation Schema}

The \texttt{annotations.json} file contains 18 Q\&A pairs (6 per document).
Each entry records: document filename, natural-language question string, exact
answer string (verbatim from the image), page number, and the bounding-box
coordinates of the answer span in normalized $[0, 1000]$ format. The 18 pairs
cover: numeric extraction (totals, subtotals), alphanumeric identifier
extraction (invoice numbers, student IDs), categorical extraction (program
names, payment terms), and multi-token phrase extraction (company names).

\subsection{Preprocessing Pipeline}

All documents were preprocessed as follows:
(1)~\textbf{PII anonymization}: genuine personal and financial identifiers
replaced with fictional substitutes (\emph{Acme Corporation},
\emph{GT-98432});
(2)~\textbf{Resolution standardization}: all images resampled to
$800{\times}1000$\,px via Lanczos downsampling;
(3)~\textbf{Adaptive threshold binarization}: receipt scans processed with
OpenCV \texttt{adaptiveThreshold} (block size 11, constant 2) to suppress
shadow gradients;
(4)~\textbf{Color normalization}: all images converted to RGB for
\texttt{UdopProcessor} compatibility.

% ═════════════════════════════════════════════════════════════════════════════
\section{Experiments and Results}
\label{sec:experiments}
% ═════════════════════════════════════════════════════════════════════════════

\subsection{Experimental Setup}

All benchmarks were conducted on Ubuntu 22.04 LTS with an Intel Core
i7-10750H (6 cores, 2.60\,GHz), 16\,GB DDR4-2933 RAM, and no dedicated GPU
(CPU-only PyTorch inference). This represents a conservative lower bound;
NVIDIA T4 GPU inference would reduce UDOP forward-pass time by approximately
$3\text{--}4\times$. FastAPI ran in single-worker mode to isolate per-request
latency.

\subsection{Latency Benchmarks}

Table~\ref{tab:latency} reports measured per-query latency across the three
cache states.

\begin{table}[t]
\caption{Query Latency by Cache State (CPU, Single Worker)}
\label{tab:latency}
\begin{center}
\renewcommand{\arraystretch}{1.2}
\begin{tabular}{|l|c|l|}
\hline
\textbf{Query Type} & \textbf{Latency} & \textbf{Operations} \\
\hline
Cold start (first query)   & ${\approx}$21.4\,s & Decode + OCR + UDOP \\
\hline
Subsequent (L2 hit)        & ${\approx}$3.8\,s  & UDOP only           \\
\hline
Repeated (L1 hit)          & 0.0\,s             & RAM lookup only     \\
\hline
\multicolumn{3}{l}{\footnotesize Speedup (L2 vs.\ cold): $\mathbf{5.6\times}$} \\
\end{tabular}
\end{center}
\end{table}

The cold-start cost decomposes as: pdf2image rendering
${\approx}$4.2\,s (3-page document at 200\,DPI) plus Tesseract OCR
${\approx}$5.1\,s/page (${\approx}$15.3\,s total) plus UDOP inference
${\approx}$1.9\,s. With the Level-2 cache populated, only UDOP inference
remains, yielding the observed $5.6\times$ speedup entirely attributable to
eliminated image decoding and OCR computation.

\subsection{Extraction Accuracy}

Table~\ref{tab:accuracy} reports representative extraction results on the
DocChat benchmark. DocChat achieves \textbf{100\% exact-match precision}
across all 18 annotated pairs.

\begin{table}[t]
\caption{Extraction Accuracy on DocChat Benchmark (Smart Mode)}
\label{tab:accuracy}
\begin{center}
\renewcommand{\arraystretch}{1.2}
\begin{tabular}{|l|l|c|}
\hline
\textbf{Document} & \textbf{Question} & \textbf{Correct?} \\
\hline
Invoice  & What is the total amount due?    & \checkmark (\$205.70) \\
Invoice  & What is the invoice number?      & \checkmark (GT-98432) \\
Invoice  & What is the GST rate?            & \checkmark (10\%)     \\
Receipt  & What is the total payment?       & \checkmark (\$47.35)  \\
Receipt  & What is the subtotal before tax? & \checkmark (\$43.50)  \\
Form     & What is the student's program?   & \checkmark (BS AI)    \\
\hline
\multicolumn{2}{|l|}{\textbf{Overall (18/18)}} & \textbf{100\%} \\
\hline
\end{tabular}
\end{center}
\end{table}

The invoice result is particularly noteworthy: \udop{} correctly isolates
\texttt{\$205.70} as the \emph{Total Amount Due} despite six other
dollar-valued fields in the OCR transcript. This demonstrates that the
spatial embedding \eqref{eq:spatial} enables the model to associate the
``Total'' row label with its right-aligned co-row value through 2D position
reasoning rather than token proximity.

\subsection{Analysis Strategy Comparison}

Table~\ref{tab:strategy} compares the three analysis strategies.

\begin{table}[t]
\caption{Analysis Strategy: Accuracy vs.\ Latency (L2-Hit State)}
\label{tab:strategy}
\begin{center}
\renewcommand{\arraystretch}{1.2}
\begin{tabular}{|l|c|c|}
\hline
\textbf{Strategy} & \textbf{Accuracy (18 pairs)} & \textbf{Avg.\ Latency} \\
\hline
Single-Page       & 72\% (13/18) & ${\approx}$3.8\,s \\
Smart (top-$N{=}2$) & \textbf{100\%} (18/18) & ${\approx}$3.8\,s \\
All-Pages           & \textbf{100\%} (18/18) & ${\approx}$7.1\,s \\
\hline
\end{tabular}
\end{center}
\end{table}

Single-Page mode fails on multi-page documents where relevant fields appear on
page two. Smart mode matches All-Pages accuracy at ${\approx}1.87\times$ lower
latency, validating the scoring formula \eqref{eq:score} for numerical and
categorical queries alike.

% ═════════════════════════════════════════════════════════════════════════════
\section{Discussion}
\label{sec:discussion}
% ═════════════════════════════════════════════════════════════════════════════

\subsection{Effect of Layout Encoding}

The 100\% benchmark precision validates the central hypothesis: trimodal
input embeddings \eqref{eq:embed} enable UDOP to correctly associate
column headers with their values across tabular row and column boundaries,
a task that defeats purely text-based models operating on linearized OCR
transcripts. The spatial coordinate embeddings encode vertical ordering
(larger $y$ values for lower rows), horizontal alignment (similar $x_0$
for left-aligned labels), and co-row membership (shared $y_0$, $y_1$
ranges), collectively resolving the multi-valued disambiguation problem.

\subsection{Limitations}

\textbf{GPU memory footprint.}
UDOP-Large requires ${\approx}3$\,GB VRAM at float32 precision. Multi-worker
deployments scale this linearly per worker; mixed-precision (float16/bfloat16)
reduces it to ${\approx}1.5$\,GB with negligible accuracy degradation.

\textbf{CPU-bound Tesseract dependency.}
Every new document incurs the full cold-start OCR cost. Tesseract accuracy
degrades with low contrast, skew, or handwriting: a $3^\circ$-rotated receipt
in our benchmark introduced word-boundary errors in ${\approx}$8\% of OCR
tokens, though UDOP's generative decoder proved robust to this level of noise.

\textbf{512-token context limit.}
Dense pages (legal contracts, multi-column layouts) may contain 600--1200 OCR
words. The current implementation truncates to the first 512 tokens
(top-left reading order), risking missed answers in lower or right-column
content. Sliding-window or hierarchical encoding strategies could address this.

\subsection{Future Work}

\textbf{Lightweight distillation.}
LayoutLMv3-base (125M parameters) achieves 82.8\% ANLS vs.\ \udop's 84.7\%---a
1.9\% accuracy trade-off for a $6\times$ parameter reduction and corresponding
inference speedup.

\textbf{Retrieval-Augmented Generation.}
For multi-document binder analysis (e.g., 500-page due-diligence data rooms),
a hybrid vector database (Pinecone, Weaviate, or pgvector) storing dense
page embeddings would enable approximate nearest-neighbor retrieval to
scale Smart-mode ranking to arbitrary corpus sizes.

\textbf{OCR-free architectures.}
Donut \cite{b_kim2022} and Nougat \cite{b_blecher2023} demonstrate
end-to-end image-to-text performance without Tesseract, eliminating both the
CPU-bound cold-start bottleneck and OCR error propagation. Integrating either
as a drop-in replacement is the most impactful single improvement for DocChat.

% ═════════════════════════════════════════════════════════════════════════════
\section{Conclusion}
\label{sec:conclusion}
% ═════════════════════════════════════════════════════════════════════════════

We presented DocChat, a fully operational Visual Document Question Answering
system that successfully operationalizes Microsoft's \udop{} transformer for
real-time interactive document information extraction. Three contributions
jointly address the deployment challenge: (1) UDOP's trimodal spatial
embeddings achieve 100\% extraction precision on a custom three-document
benchmark by correctly resolving layout-dependent field associations that
defeat all text-only baselines; (2) a dual-level MD5-fingerprinted caching
architecture reduces cold-start latency by $5.6\times$, enabling interactive
query throughput; and (3) a TF-weighted Smart Page Ranking algorithm matches
exhaustive All-Pages accuracy at roughly half the inference cost. Together,
these contributions validate a general engineering principle: layout-aware
multimodal pretraining, paired with intelligent amortization of preprocessing
costs, can bridge the gap between computationally heavy academic vision-language
models and the latency requirements of commercial real-time applications.
Future directions include lightweight model distillation, RAG-based
multi-document retrieval, and migration to fully OCR-free end-to-end
architectures such as Donut or Nougat.

% ─── ACKNOWLEDGMENT ──────────────────────────────────────────────────────────
\section*{Acknowledgment}

The authors thank the FAST School of Computing, NUCES Islamabad, for
providing the computational resources and academic environment supporting this
work, and the Computer Vision (Spring 2026) course instructors for their
guidance throughout the project.

% ─── REFERENCES ──────────────────────────────────────────────────────────────
\begin{thebibliography}{00}

\bibitem{b_mathew2021}
M.~Mathew, D.~Karatzas, and C.~V.~Jawahar,
``DocVQA: A dataset for VQA on document images,''
in \emph{Proc. IEEE/CVF Winter Conf. Applications of Computer Vision (WACV)},
2021, pp.\ 2200--2209.

\bibitem{b_tang2023}
J.~Tang, C.~Yang, S.~Wang, H.~Xu, B.~Gao, C.~Shi \emph{et al.},
``Unifying vision, text, and layout for universal document processing,''
in \emph{Proc. IEEE/CVF Conf. Computer Vision and Pattern Recognition (CVPR)},
2023, pp.\ 19254--19264.

\bibitem{b_xu2020}
Y.~Xu, M.~Li, L.~Cui, S.~Huang, F.~Wei, and M.~Zhou,
``LayoutLM: Pre-training of text and layout for document image understanding,''
in \emph{Proc. ACM SIGKDD Int. Conf. Knowledge Discovery \& Data Mining},
2020, pp.\ 1192--1200.

\bibitem{b_xu2021}
Y.~Xu, Y.~Xu, T.~Lv, L.~Cui, F.~Wei, G.~Wang \emph{et al.},
``LayoutLMv2: Multi-modal pre-training for visually-rich document understanding,''
in \emph{Proc. ACL-IJCNLP},
2021, pp.\ 2579--2591.

\bibitem{b_huang2022}
Y.~Huang, T.~Lv, L.~Cui, Y.~Lu, and F.~Wei,
``LayoutLMv3: Pre-training for document AI with unified text and image masking,''
in \emph{Proc. ACM Int. Conf. Multimedia (MM)},
2022, pp.\ 4083--4091.

\bibitem{b_kim2022}
G.~Kim, T.~Hong, M.~Yim, J.~Park, J.~Yim, W.~Hwang \emph{et al.},
``OCR-free document understanding transformer,''
in \emph{Proc. European Conf. Computer Vision (ECCV)},
2022, pp.\ 498--517.

\bibitem{b_blecher2023}
L.~Blecher, G.~Cucurull, T.~Scialom, and R.~Stojnic,
``Nougat: Neural optical understanding for academic documents,''
\emph{arXiv preprint arXiv:2308.13418}, 2023.

\bibitem{b_smith2007}
R.~Smith,
``An overview of the Tesseract OCR engine,''
in \emph{Proc. 9th Int. Conf. Document Analysis and Recognition (ICDAR)},
2007, vol.~2, pp.\ 629--633.

\bibitem{b_wolf2020}
T.~Wolf, L.~Debut, V.~Sanh, J.~Chaumond, C.~Delangue, A.~Moi \emph{et al.},
``Transformers: State-of-the-art natural language processing,''
in \emph{Proc. EMNLP System Demonstrations},
2020, pp.\ 38--45.

\bibitem{b_mathew2022}
M.~Mathew, V.~Bagal, R.~Tito, D.~Karatzas, E.~Valveny, and C.~V.~Jawahar,
``InfographicVQA,''
in \emph{Proc. IEEE/CVF Winter Conf. Applications of Computer Vision (WACV)},
2022, pp.\ 1697--1706.

\bibitem{b_msudop}
Microsoft Research,
``UDOP: Unified document processing model card,''
Hugging Face Hub, 2022.
[Online]. Available: \url{https://huggingface.co/microsoft/udop-large-512-300k}

\bibitem{b_vaswani2017}
A.~Vaswani, N.~Shazeer, N.~Parmar, J.~Uszkoreit, L.~Jones, A.~N.~Gomez,
\L.~Kaiser, and I.~Polosukhin,
``Attention is all you need,''
in \emph{Advances in Neural Information Processing Systems (NeurIPS)},
vol.~30, 2017, pp.\ 5998--6008.

\end{thebibliography}

\end{document}